\section{Conclusion}
\label{sec:conclusion}

This paper establishes three findings. First, payment-friction increases reduce trade in complex products differentially. FATF greylisting---a country-level designation issued after a multi-year AML/CFT compliance review and unrelated to bilateral goods trade---generates a 10\% additional trade penalty per standard deviation of product complexity in the panel pair-FE specification, and yearly cross-sectional estimates lie between $-0.12$ and $-0.35$ across 2018--2023. The FDI-derisked corridor measure, a broader reduced form, produces a 13\% penalty in the panel and a yearly range of $-0.21$ to $-0.37$. Second, payment-friction reductions raise complex trade differentially. EU AMLD harmonization, effective in 2017, increases complex-product trade in EU corridors by 15\% in the panel and 10--15\% in the yearly cross-sections. The sign mirror---friction up hurts, friction down helps---is the cleanest causal-mechanism evidence available: same channel, opposite shocks, opposite signs. Third, the Stablecoin Opportunity Score concentrates where regulatory harmonization is absent and friction exposure is high. The composite top five are Haiti, Fiji, Syria, Lebanon, and Nicaragua. Sixty countries fall in the actionable quadrant where capital-account openness and infrastructure readiness together support deployment.

These findings carry policy implications. Stablecoins are best understood as a substitute for missing regulatory harmonization. The AMLD result implies that when jurisdictions agree on common AML and KYC standards, complex trade flows respond on a measurable margin; when they do not, the per-corridor compliance burden falls disproportionately on complex goods. Stablecoin payment rails offer an alternative settlement layer that bypasses the bilateral correspondent-banking network in which compliance frictions accumulate. The substitution is partial, not total. Stablecoins do not eliminate the need for AML/CFT compliance---greylisting and de-risking remain a constraint regardless of the settlement technology---and effective deployment requires complementary investment in regulatory infrastructure on the receiving side. Priority targets combine high SOS, friction exposure, and existing infrastructure capacity: within the actionable quadrant, Panama, Trinidad and Tobago, the United Arab Emirates, Bahrain, and Georgia top the deployment list.

Several limitations warrant acknowledgment. The AMLD coefficient admits a rerouting interpretation: complex trade may have shifted into EU corridors as non-EU corridors got greylisted. The triple horse race controls directly for FATF exposure and the AMLD coefficient survives, but a fully clean separation of harmonization from diversion would require event-study designs around major FATF designations and is left for future work. The kitchen-sink panel specification, which adds RTA and UN-voting controls, halves the sample to 1.37 million observations as coverage drops; in the restricted sample several coefficients attenuate to non-significance. We treat this attenuation as a power problem driven by coverage, not as evidence against the mechanism, but a richer panel covering more corridor-years would settle the question. The extensive-margin specification estimates a positive but insignificant complexity-specific differential ($\hat{\delta}_5 = +0.274$, $p = 0.122$); only the unconditional level effect of country-level de-risking on the probability of RCA ($\hat{\beta}_{\text{PF}} = -0.50$, $p = 0.028$) supports the diversification channel. The country-level friction proxy is too coarse to identify within-country across-product variation, and bilateral correspondent-banking data would sharpen the test. The lead-test result is inconclusive: future de-risking predicts current trade reductions ($\hat{\gamma}^{F1} = -0.355$, $p < 0.001$) at a magnitude similar to the contemporaneous and lagged estimates. We read this as corridor-level persistence in de-risking status rather than anticipation, but the persistence-vs-anticipation ambiguity cannot be fully resolved without an event-time design that exploits the timing of the FDI threshold crossing. Two of five gravity-PCI placebo interactions (contiguity, common colony) load with significant negative coefficients alongside the headline; we interpret these as institutional-tie correlates of payment infrastructure rather than as evidence against the mechanism, but a sharper falsification would isolate payment-specific frictions from broader institutional ties. Our FDI proxy stands in for direct correspondent-banking data: bilateral FDI declines are a coarse proxy for correspondent-account closures, and direct CBR data from the BIS or SWIFT would provide a cleaner measure. The SOS is a static 2023 snapshot that does not model dynamic feedback between stablecoin adoption and trade growth.

Future work should pursue several extensions. First, direct correspondent-banking data from the SWIFT network or BIS CPMI would enable cleaner identification and resolve both the FDI-versus-CBR distinction and the country-level-vs-bilateral coarseness that limits the extensive-margin test. Second, firm-level trade data would allow testing whether the complexity penalty operates through firm exit, reduced shipment frequency, or product downgrading within firms. Third, an event-time design exploiting the timing of FDI threshold crossings or FATF designations would separate the persistence and anticipation interpretations of the lead-test result, and a Method-of-Reflections complexity index would test sensitivity of the rankings to alternative complexity measurement. Fourth, replacing the GDP-per-capita placeholder for crypto on-ramp infrastructure with the Chainalysis Global Crypto Adoption Index would give the feasibility quadrant a behavioral rather than developmental second axis. Fifth, a dynamic SOS incorporating feedback loops between payment innovation and trade would improve policy targeting. Sixth, as stablecoin adoption progresses, quasi-experimental evaluation of actual deployment episodes would provide direct evidence of the trade-creation effects that the SOS predicts.
